% ═════════════════════════════════════════════════════════════════════════════
\section{Early Foundations}
% ═════════════════════════════════════════════════════════════════════════════

The questions studied in this chapter were first crystallized in a body of work
from the late 1980s and 1990s that established both the model and the basic
techniques of parameterized verification. A foundational formulation considers a
system composed of a single control process together with an unbounded
collection of user processes sharing a common definition, communicating through
rendez-vous synchronization in the style of CCS~\cite{GerSis92}. The central
contribution was to show that reasoning about the entire infinite family can be
reduced to a finite analysis: decision procedures were given for checking
whether every process satisfies a given indexed temporal specification, and the
reachability questions arising in the analysis were connected to the theory of
Petri nets (vector addition systems), whose coverability problem is decidable.
This early demarcation of the boundary between decidable and undecidable cases,
and of the algebraic structures underlying it, set the agenda for much of the
work that followed.

A complementary early line of work targeted linear networks of arbitrary size
and introduced automata as inductive invariants~\cite{SisGyu99}. A correctness
property is itself specified by an automaton, from which an invariant automaton
is synthesized; an induction on the number of processes then establishes that
the property is preserved as the network grows. The method was shown to be
complete for safety properties. It directly anticipates both the
network-invariant approach of Section~5 and the regular model checking
framework of Section~4, in which sets of global configurations of unbounded
width are represented symbolically by finite automata.

A further theme running through this early work is the exploitation of the
symmetry inherent in a population of identical processes. Because the identity
of an individual process is immaterial to the global behavior of the system,
verification can be carried out over equivalence classes of configurations
rather than over configurations themselves~\cite{EmeSis96}. Combining this
observation with fairness leads to on-the-fly model checking procedures that
exploit symmetry to keep the explored state space manageable~\cite{GyuSis97}.
Taken together, these strands, finite reductions for synchronizing populations,
automata-based invariants for linear networks, and symmetry reduction for
replicated components, form the conceptual foundation on which the remainder
of this chapter builds.
